\documentclass[11pt]{article}
\usepackage[margin=1in]{geometry}
\usepackage{amsmath,amssymb,amsthm,mathtools,booktabs,longtable}
\usepackage[T1]{fontenc}
\usepackage{lmodern,microtype,graphicx}
\emergencystretch=2em
\usepackage[hidelinks]{hyperref}
\usepackage{fancyvrb}
\newtheorem{theorem}{Theorem}[section]
\newtheorem{lemma}[theorem]{Lemma}
\newtheorem{proposition}[theorem]{Proposition}
\newtheorem{corollary}[theorem]{Corollary}
\theoremstyle{remark}\newtheorem{remark}[theorem]{Remark}
\newcommand{\R}{\mathbb R}\newcommand{\C}{\mathbb C}\newcommand{\Q}{\mathbb Q}
\newcommand{\Z}{\mathbb Z}\newcommand{\F}{\mathbb F_2}
\newcommand{\eps}{\varepsilon}\newcommand{\Gr}{\operatorname{Gr}}
\newcommand{\rank}{\operatorname{rank}}\newcommand{\ch}{\operatorname{ch}}
\newcommand{\Span}{\operatorname{span}}\newcommand{\ind}{\operatorname{ind}}
\newcommand{\sgn}{\operatorname{sgn}}\newcommand{\sq}{\operatorname{popcount}}
\newcommand{\fnref}[1]{\footnote{See source \cite{#1}.}}
\title{RA-17: Linear Recovery and Its Topological Relaxation}
\author{Sidney Holden\\{\small Center for Computational Biology}\\{\small Flatiron Institute, Simons Foundation}}\date{13 September 2026}
\begin{document}\maketitle
\noindent\textit{Submitted edition: the partial results passed a separate informal Codex AI-agent audit on 13 September 2026; see the accompanying independent-review.md. This is not external human peer review or formal verification.}

\begin{abstract}
This is a continuation of the two preceding RA-17 proof packages, not a complete
classification of the minimum number of linear measurements. The main additional
result is an exact solution of a topological relaxation: one measurement below the
complex generic count, a continuous odd nonvanishing map exists if and only if the
complex determinantal degree is even. The corresponding evaluation bundle also
admits the required number of independent sections exactly in the even-degree case.
These statements explain a genuine limit of the earlier bundle strategy, rather
than an unevaluated characteristic number. A direct rank-one index calculation
recovers the lower bound of nineteen measurements in dimension six and shows
precisely why it does not exclude nineteen. Integer measurement constructions,
rational arithmetic checks, and a separately labelled audit of the interrupted
small-pencil calculation accompany the proofs. The existence or impossibility of
a seventeen-dimensional real $6\times6$ minimum-rank-three pencil remains
undetermined here. Consequently the overall status is \textbf{PARTIAL}.
\end{abstract}
\tableofcontents

\section{The exact target and the scope of this continuation}
For integers $1\le r$ and $2r<d$, write $\mu_\R(d,r)$ for the smallest $m$ for
which a linear map
\[
\mathcal A:M_d(\R)\longrightarrow\R^m
\]
is injective on all matrices of rank at most $r$. This is the all-parameter
classification requested by RA-17.\fnref{ra17} Define
\[
 k=2r,\qquad c=d-k,\qquad m_0=d^2-c^2=4r(d-r),\qquad n=m_0-1.
\]
The dimension $n$ is odd. The complex projective variety of matrices of rank at
most $k$ has dimension $n$ and degree
\begin{equation}\label{degree}
 D_{d,r}=\prod_{i=0}^{c-1}
 \frac{(d+i)!\,i!}{(k+i)!\,(c+i)!}.
\end{equation}
This is the standard square determinantal degree; it is also used in the real
recovery literature.\fnref{xu}

The distinction between three questions is essential. The first asks for a
\emph{linear} $\mathcal A$ that does not vanish at any nonzero rank-at-most-$k$
matrix. The second asks for a continuous odd nonvanishing map on the same unit
link. The third asks whether the evaluation bundle has a specified number of
independent sections. A positive answer to the first implies the other two.
Neither converse is proved here. In particular, arbitrary sections of an
evaluation bundle need not be restrictions of a fixed space of constant matrices.

The preceding package gives, among other results, the intervals
\[
19\le\mu_\R(6,1)\le20,\qquad
35\le\mu_\R(10,1)\le36.
\]
The present argument independently checks the first lower bound in a short
specialized calculation. It does not improve these intervals. The earlier
archives are retained unchanged as provenance; retaining them does not mean that
every earlier topological argument has been independently formalized or
peer-reviewed. No claim of historical priority is made for the results below.

\section{Linear kernels and integer upper-bound constructions}
\begin{lemma}[Kernel criterion]\label{kernel}
Let $K=\ker\mathcal A$. Uniform injectivity on rank-at-most-$r$ matrices is
equivalent to
\[
 A\in K\setminus\{0\}\quad\Longrightarrow\quad\rank A\ge k+1.
\]
Consequently, if $\kappa(d,r)$ is the maximum dimension of such a real linear
space, then $\mu_\R(d,r)=d^2-\kappa(d,r)$.
\end{lemma}
\begin{proof}
Rank subadditivity puts every difference of two rank-at-most-$r$ matrices in the
rank-at-most-$2r$ set. Conversely, express a rank-$\ell$ matrix, $\ell\le2r$,
as a sum of $\ell$ rank-one matrices, split the summands into two groups of size
at most $r$, and put a minus sign into the second group. This expresses the
matrix as the difference of two admissible signals. A nonzero such matrix in the
kernel is therefore exactly a collision. Redundant measurement rows can be
removed, and a basis of the orthogonal complement of any admissible kernel
supplies $d^2-\dim K$ measurements.
\end{proof}

\begin{proposition}[Explicit integer measurements]\label{anti}
For every admissible $(d,r)$, there is an integer-valued measurement system with
$m_0=4r(d-r)$ independent rows satisfying Lemma~\ref{kernel}.
\end{proposition}
\begin{proof}
Treat each anti-diagonal separately. If its length is $\ell$, enumerate its
entries by $j=1,\ldots,\ell$ and measure their moments
\[
 \sum_{j=1}^{\ell}j^a X_j,
 \qquad a=0,\ldots,\min(k,\ell)-1.
\]
The Vandermonde matrix has full row rank. If a vector in its kernel had at most
$k$ nonzero coordinates, the corresponding Vandermonde columns would be
independent, so that vector would be zero. Thus a nonzero kernel vector on an
anti-diagonal has at least $k+1$ nonzero entries.

Take the first nonzero anti-diagonal of a nonzero full kernel matrix and select
$k+1$ nonzero positions. In increasing row order and decreasing column order,
the resulting square submatrix is triangular: entries on one side of its
selected diagonal belong to earlier, zero anti-diagonals. Its determinant is,
up to sign, the product of the selected nonzero entries. Hence the matrix has
rank at least $k+1$. Finally,
\[
 \dim K=\sum_{\text{anti-diagonals}}\max(\ell-k,0)=(d-k)^2=c^2.
\]
The measurements have disjoint anti-diagonal supports, so their full row rank
follows from the rank of the individual Vandermonde blocks.
\end{proof}

The exported file \path{data/antidiagonal_6_1.json} contains twenty explicit
integer measurement rows and a rational basis of their sixteen-dimensional
kernel. This is the certified upper endpoint, not a claimed nineteen-measurement
construction. Other exported examples use exactly the same proof.

\begin{proposition}[Odd degree gives the exact generic count]\label{oddexact}
If $D_{d,r}$ is odd, then $\mu_\R(d,r)=m_0$.
\end{proposition}
\begin{proof}
Suppose a kernel of dimension at least $c^2+1$ avoided the real rank-at-most-$k$
cone. Restrict to a $(c^2+1)$-dimensional subspace. Avoidance is open on the real
Grassmannian: the unit rank-at-most-$k$ link is compact, and its distance from
the projectivized kernel is positive. A sufficiently small generic real
perturbation would therefore still avoid real points. But its complex
projectivization meets the complex determinantal variety transversely in
$D_{d,r}$ points. Nonreal points occur in conjugate pairs; odd degree forces a
real point. This is a contradiction. Proposition~\ref{anti} supplies equality.
\end{proof}

\section{An exact solution of the continuous odd relaxation}
Let
\[
 X=\{[A]\in\mathbb{RP}^{d^2-1}:\rank A\le k\},\qquad
 Z=\{A:\rank A\le k,\ \|A\|_F=1\}.
\]
Let $\lambda$ be the tautological real line bundle on $X$, and write
$a=w_1(\lambda)$. The map $Z\to X$ is its two-sheeted sphere bundle.

\begin{theorem}[Critical continuous relaxation]\label{continuous}
There exists a continuous map
\[
 f:Z\longrightarrow S^{n-1},\qquad f(-A)=-f(A),
\]
if and only if $D_{d,r}$ is even. Equivalently, the bundle $n\lambda$ over $X$
has a nowhere-zero section if and only if $D_{d,r}$ is even.
\end{theorem}
This assertion is at the critical count $n=m_0-1$. It does not purport to find
the smallest count for the continuous relaxation, and it does not assert that
$f$ is the normalization of a linear map.

\subsection{Orientability and the top local-coefficient group}
\begin{lemma}\label{orient}
The smooth rank-$k$ stratum $U\subset X$ is connected and orientable. The
restriction of $\lambda$ to $U$ has nontrivial sign monodromy. If
$\Z_\lambda$ denotes the corresponding sign local system, then
\[
 H^n(X;\Z_\lambda)\cong\Z/2,
 \qquad H^n(X;\F)\cong\F,
\]
and reduction modulo two between these top cohomology groups is an isomorphism.
\end{lemma}
\begin{proof}
Let $G=\Gr_k(\R^d)$ and let $S$ be its tautological rank-$k$ bundle. The
projective bundle
\[
 \pi:P(dS)\longrightarrow G
\]
resolves $X$ by recording the image plane of a matrix. It is a diffeomorphism
over $U$. Put $u=w_1(S)$ and $h=w_1(\lambda)$ on this projective bundle. The
standard tangent-bundle formulas give
\[
 w_1(TG)=d u,\quad w_1(dS)=d u,\quad
 w_1\bigl(TP(dS)\bigr)=dk\,h+\pi^*(d u+d u)=0.
\]
Here $dk$ is even because $k=2r$. Thus the resolution, and its open subset
$U$, are orientable. These formulas follow from
$TG\cong\operatorname{Hom}(S,S^\perp)$ and
$T_{\mathrm{vert}}P(E)\oplus\eps\cong\lambda^*\otimes\pi^*E$.\fnref{ms}

The full-rank $k\times d$ matrices form a connected space when $k<d$:
orthonormal row frames are connected in this range, and the positive-definite
factor in a polar decomposition is connected. Together with connectedness of
$G$, this proves connectedness of $U$. Since $k$ is even, $-I_k\in SO(k)$ can
be joined to $I_k$ by a path. Applying that path on the image of a fixed
rank-$k$ matrix joins $A$ to $-A$ without reducing the rank. Projectivization
produces a loop on which $\lambda$ reverses sign.

The singular subset $Y=X\setminus U$ has codimension
\[
 \dim X-\dim Y=2(d-k)+1=2c+1\ge3.
\]
Compact semialgebraic sets admit finite triangulations.\fnref{bcr} The long
exact sequence for $(X,Y)$ therefore identifies top cohomology of $X$ with
compactly supported top cohomology of $U$, with either coefficient system.
Poincare duality on the connected orientable manifold $U$ gives
\[
 H_c^n(U;\Z_\lambda)\cong H_0(U;\Z_\lambda)=\Z/2.
\]
The last equality is the coinvariant calculation: transport around the loop
above imposes $v=-v$. With mod-two coefficients the result is $\F$.
The coefficient sequence
$0\to\Z_\lambda\xrightarrow{2}\Z_\lambda\to\F\to0$ now shows that
reduction in degree $n$ is an isomorphism; multiplication by two is zero and
there is no cohomology above dimension $n$.
\end{proof}

\subsection{The Euler obstruction and degree parity}
\begin{proof}[Proof of Theorem~\ref{continuous}]
A nowhere-zero section of the rank-$n$ bundle $n\lambda$ is a section of its
sphere bundle with fiber $S^{n-1}$. Since $X$ is an $n$-dimensional finite CW
space and $S^{n-1}$ is $(n-2)$-connected, its only obstruction is the Euler
class in degree $n$. This is obstruction theory for vector-bundle sections,
with the orientation local system included.\fnref{hatcher}

Because $n$ is odd, that orientation system is $\Z_\lambda$. The Euler class
reduces to
\[
 w_n(n\lambda)=a^n.
\]
By Lemma~\ref{orient}, the Euler class is zero exactly when $a^n$ is zero.
Intersecting $X$ with $n$ generic real projective hyperplanes computes the
pairing of $a^n$ with its mod-two fundamental class. The intersection avoids
$Y$ and is transverse. Its complexification consists of $D_{d,r}$ points,
so its real cardinality is congruent to $D_{d,r}$ modulo two. Since top
mod-two cohomology is one-dimensional, $a^n=0$ exactly when $D_{d,r}$ is even.

Finally a section of $n\lambda$ pulls back to $Z$ in the form
$f(A)\otimes A$. Independence of the choice of $A$ or $-A$ is precisely
$f(-A)=-f(A)$. A nowhere-zero section can be normalized to have $\|f(A)\|=1$.
This proves both equivalences.
\end{proof}

\begin{corollary}[Polynomial, but not necessarily linear, realization]\label{poly}
When $D_{d,r}$ is even, there exist $n$ rational homogeneous polynomials of one
common odd degree with no common nonzero zero on the rank-at-most-$k$ cone.
No bound of one on that degree is asserted.
\end{corollary}
\begin{proof}
Approximate the coordinates of the map in Theorem~\ref{continuous} uniformly
on compact $Z$ by polynomial restrictions. Polynomials separate points, so
Stone--Weierstrass applies. Replace a polynomial $p(A)$ by
$(p(A)-p(-A))/2$ to make it odd without losing uniform accuracy. Approximate
the finitely many coefficients by rationals. Accuracy less than half the
original nonvanishing margin preserves nonvanishing.

All surviving monomials have odd total degree. Choose a common odd degree
$q$ at least as large as every degree present. Multiply each monomial of degree
$j$ by $(\sum A_{uv}^2)^{(q-j)/2}$. This leaves its value unchanged on $Z$
and makes the whole family homogeneous of degree $q$. Homogeneity extends
nonvanishing from $Z$ to the nonzero cone.
\end{proof}

The distinction from RA-17 is not cosmetic. These polynomials need not be
linear, and in general $P(X)-P(Y)\ne P(X-Y)$. The kernel/difference argument
in Lemma~\ref{kernel} cannot be transferred to such a nonlinear family.

\section{The evaluation-bundle relaxation has the same critical criterion}
Put $B=\Gr_c(\R^d)$, let $Q$ be its tautological rank-$c$ bundle, and put
\[
 E=dQ,\qquad N=dc,\qquad b=\dim B=kc,
 \qquad q=c^2+1=N-b+1.
\]
An admissible $q$-dimensional matrix space would give $q$ independent sections
of $E$ by restricting matrices to the planes represented by $Q$. If a
restriction vanished, the matrix would factor through a $k$-dimensional
quotient and have rank at most $k$.

\begin{theorem}[Exact critical frame criterion]\label{frames}
The bundle $E=dQ$ has $c^2+1$ pointwise independent continuous sections if and
only if $D_{d,r}$ is even. Equivalently, in the even-degree case there exists a
real bundle $\eta$ of rank $b-1$ such that
\[
 \eps^{c^2+1}\oplus\eta\cong dQ.
\]
\end{theorem}
\begin{proof}
Here $b=kc$ is even and $q\ge2$. The Stiefel manifold $V_{N,q}$ is
$(b-2)$-connected, and its first nonzero homotopy group is
\[
 \pi_{b-1}(V_{N,q})\cong\Z/2.
\]
For completeness, fix the first $q-2$ frame vectors. The fiber is
$V_{b+1,2}$, and the base $V_{N,q-2}$ is $b$-connected, so the relevant
homotopy group is unchanged. The two-frame manifold $V_{b+1,2}$ is the unit
tangent sphere bundle of $S^b$. In its homotopy sequence, the boundary
$\pi_b(S^b)\to\pi_{b-1}(S^{b-1})$ is multiplication by the Euler number
$\chi(S^b)=2$. Its cokernel is $\Z/2$. The case $q=2$ uses this last
calculation directly.

Since the base has dimension $b$, there is only this first obstruction to a
$q$-frame. It is the class $w_b(E)$, with untwisted mod-two coefficients.
There are no higher obstructions on this base. This is the usual
obstruction-theoretic definition of the relevant Stiefel--Whitney class.

To evaluate $w_b(E)$, take $q$ generic constant matrices and their restriction
sections. A linear dependence at a $c$-plane means that a nonzero member of
the matrix pencil annihilates that plane. Generically such a member has rank
exactly $k$ and has a unique $c$-dimensional nullspace. The degeneracy points
therefore correspond exactly to the finite projective intersection with the
rank-at-most-$k$ variety. The parity of their count is $D_{d,r}$. Top
mod-two cohomology of the connected closed manifold $B$ is $\F$, proving
that $w_b(E)$ vanishes exactly in the even-degree case. A frame has an
orthogonal complement, giving the asserted bundle decomposition.
\end{proof}

\begin{corollary}[A precise limit of bundle-only arguments]\label{limit}
When $D_{d,r}$ is even, no contradiction to the abstract decomposition
\[
 \eps^{c^2+1}\oplus\eta\cong dQ
\]
can rule out the critical kernel dimension $c^2+1$: such a decomposition
actually exists. In particular, characteristic classes or indices derived
only from that decomposition cannot exclude it.
\end{corollary}
This does not say that all topological arguments about the original linear
problem are exhausted. An argument using extra compatibility of constant
matrix sections, or extra geometry beyond this bundle relation, would be a
different argument. It says exactly which relaxation has already become
feasible.

\section{A direct rank-one index calculation}
This section specializes and independently rederives part of the earlier
index obstruction. It uses the same kernel-to-bundle principle, not a new
construction. Write $d=2\ell\ge4$, $r=1$, and set
\[
 h=\sq(\ell-1),\qquad c=d-2.
\]

\begin{theorem}[Rank-one index bound]\label{rankone}
For even $d\ge4$,
\begin{equation}\label{rankonebound}
 \mu_\R(d,1)\ \ge\ 4d-4-2h+(h\bmod2).
\end{equation}
In particular, for $d=2^s+2$ with $s\ge1$, the bound is $4d-5$.
\end{theorem}

\subsection{The relevant characteristic number}
Work over $G=\Gr_2(\R^d)$, with tautological $S$ and complementary $Q$.
An admissible kernel of dimension $\kappa=c^2+j$ forces
\[
 \eps^{\kappa}\oplus\eta\cong dQ,
 \qquad t:=\rank\eta=2(d-2)-j.
\]
The rational cohomology ring is
\[
 H^*(G;\Q)=\Q[p]/(p^\ell),\qquad |p|=4,
 \qquad p=p_1(S),\qquad \int_G p^{\ell-1}=1
\]
for a suitable orientation. The rational ring is the even/even unoriented
Grassmannian presentation.\fnref{carlson} For the normalization, the oriented
two-plane cover identifies with the complex quadric of dimension $d-2$;
the Euler class of the oriented two-plane bundle is, up to sign, its
hyperplane class. The top hyperplane number of a quadric is two, and the
cover has degree two.

Both $TG$ and $\eta$ are orientable and admit $\operatorname{Spin}^c$
structures whose determinant classes are torsion. Indeed, if
$u=w_1(S)=w_1(Q)$, then
\[
 w_2(TG)=(\ell-1)u^2,
 \qquad w_2(\eta)=\ell u^2.
\]
The integral Bockstein $\beta(u)$ reduces modulo two to $u^2$ and is
2-torsion, so it supplies the required integral lifts. Thus the
$\operatorname{Spin}^c$ index theorem applies, with rational determinant
factors equal to one.\fnref{lm}

Introduce the formal Euler root $z$ of $S$, so $p=z^2$. Stable identities give
\[
 TG\oplus(S\otimes S)\cong dS,\qquad Q\oplus S\cong\eps^d.
\]
The normalized spinor Chern character of $\eta$ is therefore
$\cosh(z/2)^{-d}$, while
\[
 \widehat A(TG)=
 \left(\frac{z/2}{\sinh(z/2)}\right)^d\frac{\sinh z}{z}.
\]
Multiplying gives $(z/\sinh z)^{d-1}$. The absolute value of the total
spinor index is consequently
\begin{equation}\label{indexformula}
 |I|=2^{\lfloor t/2\rfloor-(d-2)}
 \binom{d-2}{(d-2)/2}
 =\frac{\binom{d-2}{(d-2)/2}}{2^{\lceil j/2\rceil}}.
\end{equation}
No assumption that $G$ is spin, rather than merely $\operatorname{Spin}^c$,
is needed in this computation.

Here is a direct coefficient verification. Formal change of variables
$v=\sinh z$ yields
\begin{align*}
 [z^{d-2}]\left(\frac z{\sinh z}\right)^{d-1}
 &=\operatorname*{Res}_{z=0}\frac{dz}{(\sinh z)^{d-1}}\\
 &=[v^{d-2}](1+v^2)^{-1/2}
 =(-1)^{(d-2)/2}\frac{\binom{d-2}{(d-2)/2}}{2^{d-2}}.
\end{align*}
This derivation avoids an unevaluated localization sum.

\subsection{Integrality and the half-spin refinement}
The 2-adic valuation of the central binomial coefficient in
\eqref{indexformula} is $h=\sq(\ell-1)$, by the factorial valuation formula.
If $h$ is even, taking $j=2h+1$ makes the total index a noninteger. Such a
kernel is impossible.

If $h$ is odd, instead take $j=2h$. Then $t\equiv2\pmod4$. The Euler class
of $\eta$ vanishes rationally because $H^*(G;\Q)$ is concentrated in degrees
divisible by four. The two half-spinor Chern characters therefore coincide
rationally. Each half-index must equal $I/2$, but
\[
 |I/2|=\frac{\binom{d-2}{(d-2)/2}}{2^{h+1}}
\]
is a noninteger. Again the proposed kernel is impossible. A larger admissible
kernel would contain the forbidden smaller-dimensional one. The maximum
possible excess $j$ is therefore at most $2h-(h\bmod2)$, proving
Theorem~\ref{rankone}.

\subsection{The six-dimensional gap, explicitly}
For $d=6$, $G=\Gr_2(\R^6)$ has
\[
 H^*(G;\Q)=\Q[p]/(p^3),\qquad \int_Gp^2=1.
\]
A hypothetical eighteen-dimensional minimum-rank-three kernel would yield
an oriented rank-six $\eta$ with
\[
 p_1(\eta)=-6p,\qquad p_2(\eta)=21p^2.
\]
One also has $p_1(TG)=2p$, $p_2(TG)=7p^2$, and hence
\[
 \widehat A(TG)=1-\frac{p}{12}.
\]
The rational Euler class of $\eta$ is zero. Its half-spinor character is
\[
 \ch(\Delta^+\eta)=4-3p+\frac54p^2.
\]
The twisted Dirac index would therefore be
\[
 \int_G\left(1-\frac p{12}\right)
 \left(4-3p+\frac54p^2\right)=\frac32,
\]
which is impossible. This proves $\kappa(6,1)\le17$ and
$\mu_\R(6,1)\ge19$.

For a seventeen-dimensional kernel, the complement instead has rank seven.
The total spinor character is $8-6p+\tfrac52p^2$, giving index $3$, an
integer. More strongly,
\[
 D_{6,1}=1764
\]
is even, so Theorem~\ref{frames} proves that a rank-seven complement with
$\eps^{17}\oplus\eta\cong6Q$ actually exists. There is no missing
bundle-only nonexistence theorem at this endpoint. What remains missing is
a \emph{linear} seventeen-dimensional pencil, or an obstruction using more
than this abstract bundle decomposition.

\section{Exact algebraic certificates and the interrupted pencil}
An algebraic candidate must be distinguished from a verified construction.
The accompanying scripts use two exact acceptance mechanisms. Numerical
optimization is permitted only as a way of proposing rational data; it is
never itself an acceptance test.

\subsection{Positive Gram forms in the minor ideal}
Suppose $A(x)=\sum_{j=0}^{q-1}x_jB_j$ is a rational pencil, and $f_i(x)$
are all its $(k+1)$-minors. A certificate consists of rational polynomials
$h_i$, a vector $z(x)$ of all monomials of some positive degree, and a rational
symmetric positive-definite matrix $H$ satisfying
\begin{equation}\label{sos}
 z(x)^THz(x)=\sum_i h_i(x)f_i(x).
\end{equation}
If all $f_i(x)$ vanished at a nonzero real $x$, the right side would be zero,
but the left side would be positive: $z(x)\ne0$ because it includes pure
powers of each coordinate. Thus the pencil is admissible. Identity
\eqref{sos} is checked coefficient by coefficient over $\Q$; positive
definiteness is checked by exact rational Schur-complement pivots.

The search script first solves a semidefinite feasibility problem, then
rounds free entries to rationals and solves the linear coefficient constraints
exactly for the remaining entries. It rejects the proposal unless both exact
checks pass. An infeasible numerical solve or a timeout is labelled
inconclusive, not a counterexample to the pencil.

\subsection{Projective charts and Hermite trace forms}
An alternative verification covers projective parameter space by the disjoint
charts
\[
 x_0=\cdots=x_{i-1}=0,\qquad x_i=1,
 \qquad i=0,\ldots,q-1.
\]
This explicitly includes every projective point, so no hyperplane at infinity
is omitted. On a chart with zero-dimensional rational minor ideal $I$, let
$T=\Q[x]/I$ and form the trace pairing
\[
 H_T(u,v)=\operatorname{Tr}(M_{uv}),
\]
where $M$ denotes multiplication on the finite-dimensional algebra.
Its signature is the number of distinct real points of the chart. Indeed,
each real local factor contributes one positive direction, each conjugate
complex pair contributes one positive and one negative direction, and the
nilradical contributes only zero directions. Local multiplicities multiply
the nonzero trace forms by positive integers and do not change this signature.
Thus signature zero proves real emptiness. A unit ideal also proves emptiness;
a positive-dimensional complex ideal is reported as inconclusive by this
particular checker.

\subsection{Status of the preserved small-pencil calculation}

No new small-pencil certificate is asserted in this continuation. The earlier Xu certificate remains available in the preserved prior archive.


\section{What remains to solve RA-17}
The exact linear classification has not been obtained. In particular, neither
of the following alternatives has been established for the first retained gap:
\begin{equation}\label{missing}
 \begin{gathered}
 \text{Construct a seventeen-dimensional }K\subset M_6(\R)\\
 \text{whose every nonzero member has rank at least three,}\\
 \text{or prove that no such linear space exists.}
 \end{gathered}
\end{equation}
The first would give $\mu_\R(6,1)=19$; the second, together with
Proposition~\ref{anti}, would give $\mu_\R(6,1)=20$.

Theorems~\ref{continuous} and \ref{frames} solve specific relaxations, not
\eqref{missing}. At this endpoint they exclude a proposed strategy: merely
strengthening the characteristic-class or index calculation for the same
abstract complement cannot prove nonexistence, because the complement exists.
They do not exclude a strategy based on linearity, simultaneous compatibility
of evaluations, real algebraic geometry, or a direct certified construction.

There are other retained gaps as well. The preceding report records
$35\le\mu_\R(10,1)\le36$, and corank-one intervals at $(32,15)$ and $(64,31)$.
This continuation does not close those intervals or independently re-prove all
of their endpoints. A complete answer requires exact values for every
admissible pair, not only a larger finite table or a decision procedure for
one fixed pair. No claim of complete solution should be inferred from the
existence of a reproducible archive.

\appendix
\section{Reproduction, assurance, and source provenance}
The primary arithmetic check is
\begin{Verbatim}[fontsize=\small]
python code/verify_relaxation.py
\end{Verbatim}
It checks the factorial degree formulas, their rank-two binary valuation,
the formal-series index coefficient, the indicated mod-two Grassmannian
class, and the ranks of the exported integer measurement systems. For the
rank-two variety the closed degree formula is
\[
 D_{d,1}=\frac{(2d-3)!\,(2d-4)!}{(d-1)!^2(d-2)!^2},
 \qquad \nu_2(D_{d,1})=2\bigl(\sq(d-1)-1\bigr).
\]
Thus this degree is odd exactly when $d-1$ is a power of two. Combining with
Proposition~\ref{oddexact} gives $\mu_\R(d,1)=4d-4$ for those dimensions.
For all other dimensions, this parity computation alone is not an exact
linear classification.

The optional pencil checker is run as
\begin{Verbatim}[fontsize=\small]
python code/certify_interrupted_pencil.py --chart 0
# Repeat for charts 1, 2, 3, 4.
python code/sos_pencil_certificate.py --degree 4
# A numerical proposal is not accepted until rational verification passes.
\end{Verbatim}
Saved Gram certificates can be verified without re-running the SDP by supplying
\verb|--verify| and the certificate filename. Such verification does not need
CVXPY. The chart method recomputes the rational Groebner basis rather than
assuming that an arbitrary stored basis generates the original minor ideal.

The files under \path{logs/} record the actual executions and their exit
statuses. A missing certificate, timeout, or positive-dimensional chart is not
labelled a proof of admissibility. Arithmetic checks certify identities and
matrix computations, not the correctness of every topological inference.
The written proofs are not proof-assistant formalizations and have not been
independently peer-reviewed as part of this archive.

The earlier archives are included under \path{prior/} without modification.
The copied interrupted scripts under \path{exploratory/} are research notes,
not additional theorems. Third-party research PDFs are not redistributed.
A SHA-256 manifest records the delivered files; archive integrity is distinct
from mathematical correctness.

\begin{thebibliography}{9}
\bibitem{ra17} OpenProblemsInNLA. \emph{RA-17: Minimum linear measurements for
uniform recovery of real low-rank matrices}. Repository problem specification.
\url{https://github.com/ajt60gaibb/OpenProblemsInNLA/tree/main/randomized-and-low-rank-approximation/RA-17}.
\bibitem{xu} Zhiqiang Xu. \emph{The minimal measurement number for low-rank
matrices recovery}. arXiv:1505.07204, 2015.
\url{https://arxiv.org/abs/1505.07204}.
\bibitem{carlson} Jeffrey D. Carlson. \emph{Grassmannians and the equivariant cohomology of isotropy actions}. arXiv:1611.01175, version 2, 2023. The even/even real
Grassmannian presentation is used here; see the source audit for the title
and version metadata retained with the earlier package.
\url{https://arxiv.org/abs/1611.01175}.
\bibitem{ms} John W. Milnor and James D. Stasheff. \emph{Characteristic Classes}.
Annals of Mathematics Studies 76, Princeton University Press, 1974.
Tangent-bundle, Euler-class, and Stiefel--Whitney-class formulas.
\bibitem{hatcher} Allen Hatcher. \emph{Vector Bundles and K-Theory}.
Author's online text. Obstructions to independent sections and Stiefel manifolds.
\url{https://pi.math.cornell.edu/~hatcher/VBKT/VB.pdf}.
\bibitem{bcr} Jacek Bochnak, Michel Coste, and Marie-Francoise Roy.
\emph{Real Algebraic Geometry}. Springer, 1998. Semialgebraic triangulation.
\bibitem{lm} H. Blaine Lawson, Jr., and Marie-Louise Michelsohn.
\emph{Spin Geometry}. Princeton University Press, 1989.
Spinor characteristic classes and the $\operatorname{Spin}^c$ index theorem.
\end{thebibliography}
\end{document}
